\section{Технологическая часть}
\subsection{Средства реализации}

Для реализации программного обеспечения выбран язык программирования C++20~\cite{cpp-standard}.
Выбор языка обусловлен:
\begin{itemize}
    \item наличием опыта работы с языком;
    \item наличием стандартной библиотеки (STL), которая предоставляет контейнеры для работы с данными;
    \item возможностью использования нативных потоков;
    \item поддержкой объектно-ориентированной парадигмы;
    \item кроссплатформенностью.
\end{itemize}

Для реализации многопоточной обработки применяются средства стандартной библиотеки C++~\cite{cpp-standard}: \texttt{std::thread} для создания потоков, \texttt{std::mutex} и \texttt{std::condition\_variable} для синхронизации, \texttt{std::future} и \texttt{std::packaged\_task} для асинхронного программирования, контейнеры \texttt{std::vector}, \texttt{std::array} и умные указатели \texttt{std::shared\_ptr} для управления памятью.

Графический интерфейс реализован с использованием фреймворка Qt 6~\cite{qt6-doc}.
Применяются модули: \texttt{QtWidgets} для создания оконного интерфейса, \texttt{QtConcurrent}
для параллельного выполнения задач, \texttt{QFutureWatcher}
для отслеживания асинхронных операций, \texttt{QPixmap} и \texttt{QPainter} для работы с изображениями.
Основные компоненты интерфейса: \texttt{QMainWindow}, \texttt{QLabel}, \texttt{QPushButton}, \texttt{QFileDialog}.
Qt обеспечивает безопасную многопоточность через механизм сигналов и слотов.

Разработка выполнена в среде разработки Visual Studio Code~\cite{vscode}. Для сборки использовался CMake~\cite{cmake}.
Для модульного тестирования применяется фреймворк Google Test~\cite{gtest}.

\subsection{Реализация проведения исследования}

Время выполнения программы фиксируется с использованием библиотеки \texttt{std::chrono}, функции \texttt{high\_resolution\_clock::now()}~\cite{cpp-standard}
с точностью до миллисекунд. Измеряется чистое время рендеринга кадра без учёта инициализации сцены и сохранения результата.


\subsection{Реализация алгоритмов}

Реализация функции алгоритма генерации части кадра представлена на листинге~\ref{list:render-task} в приложении Б.

Реализация функции получения луча из камеры представлена на листинге~\ref{list:get-ray} в приложении Б.

Реализация функции получения цвета луча представлена на листинге~\ref{list:ray-color} в приложении Б.

Реализация функции пересечения с приземлённым туманом на листинге~\ref{list:hit-fog} в приложении Б.


\subsection{Пользовательский интерфейс}

Главное окно программы представлено на рисунке~\ref{fig:main_window}. Оно разделено на две части: окно визуализации и окно настроеек сцены и генерации.

\begin{figure}[H]
	\centering
	\includegraphics[width=1\textwidth]{img/3-1.png}
	\caption{главное окно программы}
	\label{fig:main_window}
\end{figure}

\subsubsection{Панель настройки объектов}

Панель настройки объектов позволяет, добавить, удалить, переместить, повернуть, изменить, сделать источником и переключить видимость объекта.
Всплывающее окно добавления объектов представлено на рисунке~\ref{fig:popup-add-object}.
Пользователю необходимо указать геометрические параметры объекта и подтвердить выбор, после чего объект будет добавлен в сцену,
и сразу же визуализирован.

\begin{figure}[H]
	\centering
	\includegraphics[width=1\textwidth]{img/3-2.png}
	\caption{окно добавление объекта}
	\label{fig:popup-add-object}
\end{figure}

\subsubsection{Панель изменения материала}

Программа позволяет поменять материал объекта, на один из заранее установленных. Окно изменения материала продемонстрировано на рисунке~\ref{fig:popup-change-mat}.
\begin{figure}[H]
	\centering
	\includegraphics[width=1\textwidth]{img/3-3.png}
	\caption{окно изменения материала}
	\label{fig:popup-change-mat}
\end{figure}

\subsubsection{Панель управления камерой}

Панель управления камерой позволяет добавить, удалить, изменить положение и сделать камеру активной.
При добавлении и изменении камеры необходимо указать позицию камеры, точку куда она смотрит, угол обзора, расстояние до точки фокуса
и угол расфокусировки. Окно добавление камеры продемонстрировано на рисунке~\ref{fig:popup-add-camera}
\begin{figure}[H]
	\centering
	\includegraphics[width=1\textwidth]{img/3-4.png}
	\caption{окно добавление камеры}
	\label{fig:popup-add-camera}
\end{figure}

\subsubsection{Панель настроек генерации}

Панель настрок генерации изображения позволяет выбрать количество лучей на 1 пиксель, а также максимальную глубину рекурсии функции трассировки пути.
При нажатии на кнопку начала генерации появляется всплывающее окно, в которое выводится результат. Процесс генерации изображения продемонстрирован на рисунке~\ref{fig:process-render}. 
\begin{figure}[H]
	\centering
	\includegraphics[width=1\textwidth]{img/3-5.png}
	\caption{процесс генерации изображения}
	\label{fig:process-render}
\end{figure}

Результат генерации изображение приведен на рисунке~\ref{fig:render-finished}
\begin{figure}[H]
	\centering
	\includegraphics[width=1\textwidth]{img/3-6.png}
	\caption{результат генерации изображения}
	\label{fig:render-finished}
\end{figure}

\subsubsection{Панель добавления тумана}

Панель позволяет добавить туман на сцену. Окно позволяет выбрать тип тумана, ограничивающий объем, плотность масштаб шума и фактор высотности.
Окно добавления тумана приведено на рисунке~\ref{fig:add-fog}
\begin{figure}[H]
	\centering
	\includegraphics[width=1\textwidth]{img/3-7.png}
	\caption{Окно добавления тумана}
	\label{fig:add-fog}
\end{figure}

\subsection*{Вывод}

В данной части был обоснован выбор средств реализации, приведены исходные коды
реализации используемых алгоритмов и приведены изображения интерфейса.

\clearpage